\documentclass[11pt]{article}

\usepackage[utf8]{inputenc}
\usepackage[T1]{fontenc}
\usepackage{lmodern}
\usepackage[margin=1in]{geometry}
\usepackage{amsmath,amssymb,amsthm}
\usepackage{booktabs}
\usepackage{array}
\usepackage[colorlinks=true,linkcolor=blue,citecolor=blue,urlcolor=blue]{hyperref}

\newcommand{\coloneqq}{\mathrel{:=}}

\setlength{\emergencystretch}{1.5em}

\theoremstyle{plain}
\newtheorem{theorem}{Theorem}[section]
\newtheorem{proposition}[theorem]{Proposition}
\newtheorem{lemma}[theorem]{Lemma}
\newtheorem{corollary}[theorem]{Corollary}

\theoremstyle{definition}
\newtheorem{definition}[theorem]{Definition}
\newtheorem{counterexample}[theorem]{Counterexample}
\newtheorem{example}[theorem]{Example}

\theoremstyle{remark}
\newtheorem{remark}[theorem]{Remark}

\newcommand{\Cl}{\mathrm{Cl}}
\newcommand{\grade}[2]{\langle #1 \rangle_{#2}}
\newcommand{\lcontract}{\mathbin{\lrcorner}}
\newcommand{\rcontract}{\mathbin{\llcorner}}

\title{\textbf{Resolution of Conjecture 00000007601:\\
Symmetric--Antisymmetric Decomposition of the Clifford Product}}
\author{Formal Mathematical Resolution \& Lean 4 Verification}
\date{\today}

\begin{document}
\maketitle

\begin{abstract}
We resolve Conjecture \texttt{00000007601} concerning the decomposition of the product in a Clifford algebra $\Cl(V, Q)$ generated by a quadratic form $Q$. The conjecture asserts that the Clifford product decomposes via the commutator structure into symmetric and antisymmetric parts, that the symmetric part is of inner-product type, that the antisymmetric part is of outer-product type, and that the orthogonality of the decomposition is governed by the metric. We prove that \textbf{the conjecture is false in general} for the Clifford algebra product (as soon as $\dim V \ge 2$), while establishing as a positive non-vacuity theorem that all clauses hold when restricted to grade-$1$ generating vectors $u, v \in V \subset \Cl(V, Q)$ (with non-isotropic metric direct-sum properties requiring nondegeneracy/definiteness of $Q$).

More precisely, we prove an exact homogeneous grade-parity sign law for $X_r \in \Cl^r(V, Q)$ and $Y_s \in \Cl^s(V, Q)$:
\[
  \grade{Y_s X_r}{r + s - 2j} = (-1)^{rs - j}\,\grade{X_r Y_s}{r + s - 2j}
  \qquad (0 \le 2j \le r + s).
\]
Consequently, for a vector $u \in \Cl^1$ and a bivector $B \in \Cl^2$, the roles asserted by the conjecture \emph{invert completely}: the symmetric anticommutator $\frac{1}{2}\{u, B\} = u \wedge B$ is the outer product, while the antisymmetric commutator $\frac{1}{2}[u, B] = u \lrcorner B$ is the inner contraction. For two bivectors in $\Cl_{4,0}(\mathbb{Q})$, the symmetric part simultaneously mixes grade-$0$ inner contraction and grade-$4$ outer product, while the antisymmetric commutator is a grade-$2$ Lie bracket belonging to neither type. Furthermore, even in Euclidean $\Cl_{2,0}(\mathbb{Q})$, inhomogeneous elements $X = 1 + e_1$ and $Y = e_2 + e_1 e_2$ yield identical nonzero symmetric and antisymmetric parts $S = A = e_2 + e_1 e_2$ with strictly positive metric inner product $\langle S, A \rangle_{\widetilde{Q}} = 2 \neq 0$. All positive theorems, general parity laws, and counterexamples are formally verified from first principles in Core Lean~4.19.0 with zero external dependencies and independently checked in exact rational arithmetic.
\end{abstract}

\section{Problem Statement and Algebraic Preliminaries}

\subsection{Statement of Conjecture 00000007601}

The source statement (\texttt{SOURCE.md}, in both English and Chinese) reads verbatim in English:
\begin{quote}
\emph{Definition: Clifford algebras as algebras generated by quadratic forms. Conjecture: The product decomposes into symmetric plus antisymmetric parts and the combination of the decomposition is the commutator structure; the symmetric part is inner-product type; the antisymmetric part is outer-product type; and the orthogonality of the decomposition is given by the metric. (symmetric-antisymmetric Clifford product decomposition)}
\end{quote}
The parallel Chinese text states identically that $\Cl(V, Q)$ is the algebra generated by a quadratic form, and conjectures four properties for the Clifford product (\emph{Clifford product symmetric--antisymmetric decomposition}):
\begin{enumerate}
  \item[\textbf{(C1)}] \textbf{Commutator--anticommutator splitting:} The Clifford product $X Y$ decomposes into symmetric plus antisymmetric parts given by the anticommutator and commutator structure.
  \item[\textbf{(C2)}] \textbf{Symmetric = inner-product type:} The symmetric part of $X Y$ is of inner-product (contraction) type.
  \item[\textbf{(C3)}] \textbf{Antisymmetric = outer-product type:} The antisymmetric part of $X Y$ is of outer-product (exterior/wedge) type.
  \item[\textbf{(C4)}] \textbf{Metric orthogonality:} The orthogonality of the decomposition is given by the metric induced by $Q$.
\end{enumerate}

Notice that the definition specifies the entire Clifford algebra $\Cl(V, Q)$ generated by $(V, Q)$, and the conjecture is formulated for \emph{the product} in $\Cl(V, Q)$ without restricting the operands to the grade-$1$ generating subspace $V \subset \Cl(V, Q)$. Below we analyze \emph{both} the restriction to grade-$1$ vectors $u, v \in V$ and the general Clifford product on homogeneous blades and general multivectors in $\Cl(V, Q)$.

\subsection{Clifford Algebras, Canonical Grading, Contractions, and Metric Forms}

Let $K$ be a field with $\mathrm{char}(K) \neq 2$ (or more generally a commutative ring in which $2$ is invertible; for concrete counterexamples we work over $K = \mathbb{Q}$). Let $V$ be an $n$-dimensional $K$-vector space equipped with a quadratic form $Q : V \to K$ and associated symmetric polar bilinear form
\begin{equation}\label{eq:polar}
  B_Q(u, v) = \frac{1}{2}\bigl(Q(u + v) - Q(u) - Q(v)\bigr), \qquad B_Q(v, v) = Q(v).
\end{equation}

\begin{definition}[Clifford Algebra]\label{def:clifford}
The \emph{Clifford algebra} $\Cl(V, Q)$ is the quotient $T(V) / I_Q$ of the tensor algebra $T(V) = \bigoplus_{k \ge 0} V^{\otimes k}$ by the two-sided ideal $I_Q$ generated by $\{v \otimes v - Q(v)\cdot 1 \mid v \in V\}$, equipped with the canonical $K$-linear embedding $\iota : V \hookrightarrow \Cl(V, Q)$ (written simply $v \in V \subset \Cl(V, Q)$). It is characterized up to unique $K$-algebra isomorphism by the universal property: for every associative unital $K$-algebra $A$ and $K$-linear map $f : V \to A$ satisfying $f(v)^2 = Q(v)\cdot 1_A$ for all $v \in V$, there exists a unique $K$-algebra homomorphism $\Phi : \Cl(V, Q) \to A$ with $\Phi \circ \iota = f$.
\end{definition}

Polarizing $(u + v)^2 = Q(u + v)\cdot 1$ yields the fundamental generator anticommutation relation
\begin{equation}\label{eq:polar-clifford}
  \{u, v\} \coloneqq u v + v u = 2 B_Q(u, v)\cdot 1 \qquad (u, v \in V).
\end{equation}
Since $\mathrm{char}(K) \neq 2$ and $\dim_K V = n < \infty$, $(V, Q)$ admits a $B_Q$-orthogonal basis $(e_1, \dots, e_n)$ with $e_i^2 = q_i \cdot 1$ ($q_i = Q(e_i)$) and $e_i e_j = -e_j e_i$ for $i \neq j$. More canonically, Chevalley's antisymmetrization map $\psi : \Lambda(V) = \bigoplus_{k=0}^n \Lambda^k(V) \xrightarrow{\sim} \Cl(V, Q)$, defined when $k! \in K^\times$ by $v_1 \wedge \dots \wedge v_k \mapsto \frac{1}{k!}\sum_{\sigma \in S_k} \mathrm{sgn}(\sigma)\,v_{\sigma(1)}\cdots v_{\sigma(k)}$ (and sending $e_{i_1} \wedge \dots \wedge e_{i_k} \mapsto e_{i_1}\cdots e_{i_k}$ for $1 \le i_1 < \dots < i_k \le n$ in any $B_Q$-orthogonal basis), is a basis-independent $K$-linear isomorphism.

\begin{definition}[Canonical Grading, Involutions, Outer/Inner Products, and Metric]\label{def:ops}
Via $\psi : \Lambda(V) \xrightarrow{\sim} \Cl(V, Q)$, we write $\Cl(V, Q) = \bigoplus_{k=0}^n \Cl^k(V, Q)$ as a graded $K$-vector space, with grade-$k$ projection $\grade{X}{k} \in \Cl^k(V, Q)$ ($\grade{X}{k} = 0$ for $k > n$).
\begin{enumerate}
  \item[(i)] \emph{Main grade involution and reversion:} The \emph{grade involution} $X \mapsto \widehat{X}$ is the algebra automorphism determined by $\widehat{v} = -v$ ($v \in V$), acting on $\Cl^k(V, Q)$ by $(-1)^k$. \emph{Reversion} $X \mapsto \widetilde{X}$ is the algebra anti-involution ($\widetilde{X Y} = \widetilde{Y}\widetilde{X}$) fixing $v \in V$, acting on $e_{i_1}\cdots e_{i_k} \in \Cl^k(V, Q)$ by
  \begin{equation}\label{eq:rev-sign}
    \widetilde{X} = (-1)^{k(k-1)/2} X \qquad \bigl(X \in \Cl^k(V, Q)\bigr).
  \end{equation}
  \item[(ii)] \emph{Symmetric and antisymmetric product parts:} For any $X, Y \in \Cl(V, Q)$,
  \begin{equation}\label{eq:sym-antisym}
    S(X, Y) \coloneqq \frac{1}{2}\{X, Y\} = \frac{1}{2}(X Y + Y X), \qquad
    A(X, Y) \coloneqq \frac{1}{2}[X, Y] = \frac{1}{2}(X Y - Y X).
  \end{equation}
  \item[(iii)] \emph{Outer product and inner contractions:} For homogeneous multivectors $X_r \in \Cl^r(V, Q)$ and $Y_s \in \Cl^s(V, Q)$, extending bilinearly to $\Cl(V, Q)$:
  \begin{align}
    X_r \wedge Y_s &\coloneqq \grade{X_r Y_s}{r + s} && \text{(exterior/outer product)},\label{eq:wedge}\\
    X_r \lcontract Y_s &\coloneqq \begin{cases}\grade{X_r Y_s}{s - r} & \text{if } r \le s,\\ 0 & \text{if } r > s,\end{cases} && \text{(left contraction)},\label{eq:lcontract}\\
    X_r \rcontract Y_s &\coloneqq \begin{cases}\grade{X_r Y_s}{r - s} & \text{if } s \le r,\\ 0 & \text{if } s > r,\end{cases} && \text{(right contraction)},\label{eq:rcontract}\\
    X_r \cdot Y_s &\coloneqq \begin{cases}\grade{X_r Y_s}{|r - s|} & \text{if } r, s \ge 1,\\ 0 & \text{if } r = 0 \text{ or } s = 0,\end{cases} && \text{(Hestenes inner product)}.\label{eq:hestenes}
  \end{align}
  \item[(iv)] \emph{Metric bilinear forms on $\Cl(V, Q)$:} Identifying $\Cl^0(V, Q) = K \cdot 1 \cong K$ via scalar extraction $\mathrm{sc}(X) \in K$ ($\grade{X}{0} = \mathrm{sc}(X)\cdot 1$), the canonical metric extension of $B_Q$ to $\Cl(V, Q)$ is the \emph{reversion scalar product}
  \begin{equation}\label{eq:metric-form}
    \langle X, Y \rangle_{\widetilde{Q}} \coloneqq \mathrm{sc}(\widetilde{X} Y),
  \end{equation}
  which satisfies $\langle u_1 \wedge \dots \wedge u_k, v_1 \wedge \dots \wedge v_k \rangle_{\widetilde{Q}} = \det\bigl(B_Q(u_a, v_b)\bigr)_{1 \le a, b \le k}$ and is positive definite whenever $Q$ is positive definite over $\mathbb{R}$ or $\mathbb{Q}$. We also consider the untwisted scalar trace $\langle X, Y \rangle_0 \coloneqq \mathrm{sc}(X Y)$, which differs from $\langle X, Y \rangle_{\widetilde{Q}}$ by $(-1)^{k(k-1)/2}$ on $\Cl^k(V, Q)$.
\end{enumerate}
\end{definition}

\section{Universal Clause (C1) and the Positive Grade-1 Theorem}

We first verify that Clause \textbf{(C1)} holds throughout $\Cl(V, Q)$, and that \emph{all four clauses} hold cleanly on the generating vector space $V = \Cl^1(V, Q)$.

\begin{proposition}[Universal Commutator--Anticommutator Decomposition]\label{prop:c1}
For all $X, Y \in \Cl(V, Q)$, we have $X Y = S(X, Y) + A(X, Y)$ with $S(X, Y) = S(Y, X)$ and $A(Y, X) = -A(X, Y)$. Moreover, $(S, A)$ is the unique pair of a symmetric operation and an antisymmetric operation on $\Cl(V, Q)$ whose sum is the Clifford product.
\end{proposition}
\begin{proof}
Adding $\frac{1}{2}(X Y + Y X)$ and $\frac{1}{2}(X Y - Y X)$ cancels $\pm \frac{1}{2} Y X$ and yields $\frac{1}{2}(2 X Y) = X Y$. Symmetry of $S$ and antisymmetry of $A$ are immediate from~\eqref{eq:sym-antisym}. Conversely, if $S'(X, Y) + A'(X, Y) = X Y$ with $S'(Y, X) = S'(X, Y)$ and $A'(Y, X) = -A'(X, Y)$, then $X Y + Y X = 2 S'(X, Y)$ and $X Y - Y X = 2 A'(X, Y)$, dividing by $2$ forces $S' = S$ and $A' = A$.
\end{proof}

\begin{theorem}[Positive Grade-$1$ Hestenes Decomposition and Lagrange--Pythagoras Identity]\label{thm:grade1}
Let $(V, Q)$ be an $n$-dimensional quadratic space over $K$ ($\mathrm{char}(K) \neq 2$). For all grade-$1$ vectors $u, v \in V = \Cl^1(V, Q)$:
\begin{enumerate}
  \item[(a)] \emph{Inner and outer identification:}
  \[
    u v = S(u, v) + A(u, v), \qquad
    S(u, v) = B_Q(u, v)\cdot 1 \in \Cl^0(V, Q), \qquad
    A(u, v) = u \wedge v \in \Cl^2(V, Q).
  \]
  When $n \ge 1$, $S(u, v) = u \lcontract v = u \rcontract v = u \cdot v$.
  \item[(b)] \emph{Metric orthogonality and Lagrange--Pythagoras identity:} Under $\langle \cdot, \cdot \rangle_{\widetilde{Q}}$, $\langle S(u, v), A(u, v) \rangle_{\widetilde{Q}} = 0$ and
  \[
    \langle u v, u v \rangle_{\widetilde{Q}} = \bigl\langle S(u, v), S(u, v) \bigr\rangle_{\widetilde{Q}} + \bigl\langle A(u, v), A(u, v) \bigr\rangle_{\widetilde{Q}} = Q(u)\,Q(v).
  \]
  \item[(c)] \emph{Anticommutation criterion:} $B_Q(u, v) = 0 \iff S(u, v) = 0 \iff u v = - v u$.
\end{enumerate}
\end{theorem}
\begin{proof}
By~\eqref{eq:polar-clifford}, $S(u, v) = \frac{1}{2}(u v + v u) = B_Q(u, v)\cdot 1 \in \Cl^0(V, Q)$. Expanding $u = \sum_i u_i e_i$ and $v = \sum_j v_j e_j$ in an orthogonal basis gives
\[
  u v = \sum_{i=1}^n q_i u_i v_i \cdot 1 + \sum_{1 \le i < j \le n} (u_i v_j - u_j v_i)\,e_i e_j = B_Q(u, v)\cdot 1 + u \wedge v,
\]
since $e_i e_j = e_i \wedge e_j \in \Cl^2(V, Q)$ for $i < j$. Subtracting $S(u, v) = B_Q(u, v)\cdot 1$ from $u v = S(u, v) + A(u, v)$ yields $A(u, v) = u \wedge v \in \Cl^2(V, Q)$, while $\grade{u v}{0} = B_Q(u, v)\cdot 1 = u \lcontract v = u \rcontract v = u \cdot v$.

For (b), distinct homogeneous grades are $\langle \cdot, \cdot \rangle_{\widetilde{Q}}$-orthogonal because for distinct multi-indices $I \neq J$, $\widetilde{e_I} e_J$ is a scalar multiple of $e_{I \triangle J}$ with $I \triangle J \neq \emptyset$, so $\mathrm{sc}(\widetilde{e_I} e_J) = 0$. Since $S(u, v) \in \Cl^0$ and $A(u, v) \in \Cl^2$, $\langle S(u, v), A(u, v) \rangle_{\widetilde{Q}} = 0$. By symmetry of $\langle \cdot, \cdot \rangle_{\widetilde{Q}}$,
\[
  \langle u v, u v \rangle_{\widetilde{Q}} = \langle S + A, S + A \rangle_{\widetilde{Q}} = \langle S, S \rangle_{\widetilde{Q}} + \langle A, A \rangle_{\widetilde{Q}}.
\]
Directly from~\eqref{eq:metric-form}, using $\widetilde{u v} = \widetilde{v}\widetilde{u} = v u$ and associativity,
\[
  \langle u v, u v \rangle_{\widetilde{Q}} = \mathrm{sc}\bigl((v u)(u v)\bigr) = \mathrm{sc}\bigl(v (u^2) v\bigr) = Q(u)\,\mathrm{sc}(v^2) = Q(u)\,Q(v).
\]
Finally, (c) follows from injectivity of $c \mapsto c\cdot 1$ and $u v + v u = 2 S(u, v)$.
\end{proof}

\begin{example}[Concrete Non-Vacuity Witness in $\Cl_{2,0}(\mathbb{Q})$]\label{ex:nonvac}
In $\Cl_{2,0}(\mathbb{Q})$ with $Q(x_1, x_2) = x_1^2 + x_2^2$, take $u = e_1 + 2 e_2$ and $v = 3 e_1 - e_2$. Then $u v = 1 - 7 e_1 e_2$ and $v u = 1 + 7 e_1 e_2$, so $S(u, v) = 1 \neq 0$ and $A(u, v) = -7 e_1 e_2 \neq 0$ are both nonzero, $\langle S, A \rangle_{\widetilde{Q}} = 0$, and $\langle u v, u v \rangle_{\widetilde{Q}} = 1^2 + (-7)^2 = 50 = 5 \cdot 10 = Q(u)\,Q(v)$.
\end{example}

\section{Exact Homogeneous Grade-Parity Law and Parity Reversal}

To understand why Clauses \textbf{(C2)}--\textbf{(C3)} fail beyond grade-$1$ vectors, we derive the exact commutation sign at every output grade for arbitrary homogeneous multivectors.

\begin{theorem}[Exact Homogeneous Grade-Parity Sign Law]\label{thm:parity}
Let $X_r \in \Cl^r(V, Q)$ and $Y_s \in \Cl^s(V, Q)$. In the Clifford product expansion
\[
  X_r Y_s = \sum_{j = 0}^{\min(r, s)} \grade{X_r Y_s}{r + s - 2j},
\]
every output grade $m = r + s - 2j$ ($0 \le 2j \le r + s$) satisfies
\begin{equation}\label{eq:parity-law}
  \grade{Y_s X_r}{r + s - 2j} = (-1)^{rs - j}\,\grade{X_r Y_s}{r + s - 2j}.
\end{equation}
Consequently, the symmetric and antisymmetric parts of $X_r Y_s$ are given by
\begin{equation}\label{eq:sym-antisym-grades}
  S(X_r, Y_s) = \sum_{\substack{0 \le j \le \min(r, s)\\ rs - j \text{ even}}} \grade{X_r Y_s}{r + s - 2j}, \qquad
  A(X_r, Y_s) = \sum_{\substack{0 \le j \le \min(r, s)\\ rs - j \text{ odd}}} \grade{X_r Y_s}{r + s - 2j}.
\end{equation}
\end{theorem}
\begin{proof}
Expanding $X_r$ and $Y_s$ in a $B_Q$-orthogonal basis $(e_1, \dots, e_n)$, if $|I| = r$ and $|J| = s$ share $|I \cap J| = j$ common indices ($0 \le j \le \min(r, s)$), then $e_I e_J$ is a scalar multiple of $e_{I \triangle J}$, which has grade $|I \triangle J| = r + s - 2j$; thus only grades $m = r + s - 2j$ can appear.

Since reversion preserves each grade subspace $\Cl^m(V, Q)$ and acts on it by $\epsilon_m \coloneqq (-1)^{m(m-1)/2}$, grade projection commutes with reversion. Applying reversion to $\grade{Y_s X_r}{m}$ in two ways gives
\[
  \epsilon_m \grade{Y_s X_r}{m}
  = \widetilde{\grade{Y_s X_r}{m}}
  = \grade{\widetilde{Y_s X_r}}{m}
  = \grade{\widetilde{X_r}\,\widetilde{Y_s}}{m}
  = \epsilon_r \epsilon_s \grade{X_r Y_s}{m}.
\]
Multiplying by $\epsilon_m$ (using $\epsilon_m^2 = 1$) yields $\grade{Y_s X_r}{m} = \epsilon_r \epsilon_s \epsilon_m \grade{X_r Y_s}{m}$. For $m = r + s - 2j$, we compute modulo $2$:
\begin{align*}
  \frac{r(r-1) + s(s-1) - m(m-1)}{2}
  &= \frac{r(r-1) + s(s-1) - (r + s - 2j)(r + s - 2j - 1)}{2}\\
  &= -rs + 2j(r + s) - 2j^2 - j \equiv rs - j \pmod{2}.
\end{align*}
Hence $\epsilon_r \epsilon_s \epsilon_m = (-1)^{rs - j}$, establishing~\eqref{eq:parity-law} and immediately implying~\eqref{eq:sym-antisym-grades}.
\end{proof}

\begin{corollary}[Universal Vector--Multivector Parity Dichotomy]\label{cor:vec-biv}
Let $u \in \Cl^1(V, Q)$ and $B_s \in \Cl^s(V, Q)$ with $1 \le s \le n$. Then $u B_s = u \lcontract B_s + u \wedge B_s$, where:
\begin{enumerate}
  \item[(i)] If $s$ is \textbf{odd}, $S(u, B_s) = u \lcontract B_s = u \cdot B_s$ and $A(u, B_s) = u \wedge B_s$.
  \item[(ii)] If $s$ is \textbf{even} (in particular for every bivector $B_2 \in \Cl^2(V, Q)$ when $n \ge 2$), the roles \textbf{reverse completely}:
  \begin{equation}\label{eq:reversal}
    S(u, B_s) = \frac{1}{2}\{u, B_s\} = u \wedge B_s, \qquad
    A(u, B_s) = \frac{1}{2}[u, B_s] = u \lcontract B_s = u \cdot B_s.
  \end{equation}
\end{enumerate}
\end{corollary}
\begin{proof}
Since $\min(1, s) = 1$, Theorem~\ref{thm:parity} allows only $j = 0$ (grade $s + 1$, equal to $u \wedge B_s$) with sign $(-1)^{s - 0} = (-1)^s$, and $j = 1$ (grade $s - 1$, equal to $u \lcontract B_s = u \cdot B_s$) with sign $(-1)^{s - 1} = -(-1)^s$. When $s$ is odd, $(-1)^s = -1$ and $(-1)^{s-1} = +1$; when $s$ is even, $(-1)^s = +1$ and $(-1)^{s-1} = -1$.
\end{proof}

\section{Explicit Counterexamples Refuting Conjecture 00000007601}

We now exhibit concrete rational counterexamples demonstrating every failure mode of Clauses \textbf{(C2)}, \textbf{(C3)}, and \textbf{(C4)}, beginning already in two-dimensional Euclidean space.

\begin{counterexample}[Vector $\times$ Bivector in $\Cl_{2,0}(\mathbb{Q})$: Antisymmetric = Inner, Symmetric = $0$]\label{cex:cl2}
Let $V = \mathbb{Q}^2$ with positive-definite Euclidean quadratic form $Q(x_1, x_2) = x_1^2 + x_2^2$ ($e_1^2 = e_2^2 = 1$, $e_1 e_2 = -e_2 e_1$). Take $A = e_1 \in \Cl^1$ and $B = e_1 e_2 \in \Cl^2$. Then:
\[
  A B = e_1(e_1 e_2) = e_2, \qquad B A = (e_1 e_2) e_1 = -e_1^2 e_2 = -e_2.
\]
Hence
\[
  S(A, B) = \frac{1}{2}\{A, B\} = 0, \qquad A_{\mathrm{anti}}(A, B) = \frac{1}{2}[A, B] = e_2 \neq 0.
\]
Meanwhile, the outer product vanishes ($A \wedge B = \grade{A B}{3} = 0$ since $\dim V = 2$), whereas the inner contractions equal $A \lcontract B = A \cdot B = \grade{A B}{1} = e_2 \neq 0$ (and $A \rcontract B = 0$, $\mathrm{sc}(A B) = 0$). Thus:
\begin{itemize}
  \item $A_{\mathrm{anti}}(A, B) = e_2 \neq 0 = A \wedge B$, refuting Clause \textbf{(C3)}; in fact $A_{\mathrm{anti}}(A, B) = A \lcontract B = A \cdot B$ is \emph{pure inner contraction}.
  \item $S(A, B) = 0 \neq e_2 = A \lcontract B = A \cdot B$, refuting Clause \textbf{(C2)} for left and Hestenes inner products.
\end{itemize}
\end{counterexample}

\begin{counterexample}[Bivector $\times$ Vector and Bivector Commutator in $\Cl_{3,0}(\mathbb{Q})$]\label{cex:cl3}
Let $V = \mathbb{Q}^3$ with $Q(x_1, x_2, x_3) = x_1^2 + x_2^2 + x_3^2$ ($e_1^2 = e_2^2 = e_3^2 = 1$).
\begin{enumerate}
  \item[(a)] Take $X = e_1 e_2 \in \Cl^2$ and $Y = e_3 \in \Cl^1$. Moving $e_3$ past both $e_1$ and $e_2$ incurs two sign changes:
  \[
    X Y = e_1 e_2 e_3, \qquad Y X = e_3 e_1 e_2 = (-1)^2 e_1 e_2 e_3 = e_1 e_2 e_3.
  \]
  Therefore $S(X, Y) = \frac{1}{2}\{X, Y\} = e_1 e_2 e_3 = X \wedge Y \neq 0$ (\emph{pure outer product}), while $A(X, Y) = \frac{1}{2}[X, Y] = 0$ and \emph{every} inner contraction vanishes ($X \lcontract Y = X \rcontract Y = X \cdot Y = \grade{X Y}{0} = 0$). Moreover, $X$ and $Y$ are mutually orthogonal under $\langle X, Y \rangle_{\widetilde{Q}} = 0$ and span $B_Q$-orthogonal subspaces $\mathrm{span}(e_1, e_2) \perp \mathrm{span}(e_3)$, yet $S(X, Y) \neq 0$ ($X$ and $Y$ commute rather than anticommute).
  \item[(b)] Take two bivectors $B_1 = e_1 e_2 \in \Cl^2$ and $B_2 = e_2 e_3 \in \Cl^2$. Then $B_1 B_2 = e_1 e_3$ and $B_2 B_1 = -e_1 e_3$, so $S(B_1, B_2) = 0$ and $A(B_1, B_2) = e_1 e_3 \neq 0 \in \Cl^2$ (the $\mathfrak{so}(3)$ Lie bracket), even though \emph{both} the outer product $B_1 \wedge B_2 = \grade{B_1 B_2}{4} = 0$ and all inner contractions ($B_1 \lcontract B_2 = B_1 \rcontract B_2 = B_1 \cdot B_2 = \grade{B_1 B_2}{0} = 0$) vanish identically.
\end{enumerate}
\end{counterexample}

\begin{counterexample}[Three-Part Bivector Product in $\Cl_{4,0}(\mathbb{Q})$]\label{cex:cl4}
Let $V = \mathbb{Q}^4$ with $Q(x) = x_1^2 + x_2^2 + x_3^2 + x_4^2$. Consider the two homogeneous bivectors
\[
  B_1 = e_1 e_2 \in \Cl^2, \qquad B_2 = e_1 e_2 + e_2 e_3 + e_3 e_4 \in \Cl^2.
\]
Using $(e_1 e_2)^2 = -1$, $(e_1 e_2)(e_2 e_3) = e_1 e_3 = -(e_2 e_3)(e_1 e_2)$, and $(e_1 e_2)(e_3 e_4) = e_1 e_2 e_3 e_4 = (e_3 e_4)(e_1 e_2)$:
\[
  B_1 B_2 = -1 + e_1 e_3 + e_1 e_2 e_3 e_4, \qquad
  B_2 B_1 = -1 - e_1 e_3 + e_1 e_2 e_3 e_4.
\]
Thus $B_1 B_2$ has \emph{three} simultaneously nonzero homogeneous components in grades $0$, $2$, and $4$:
\begin{align*}
  S(B_1, B_2) &= \frac{1}{2}\{B_1, B_2\} = \underbrace{-1}_{=\, B_1 \cdot B_2 \text{ (inner)}} + \underbrace{e_1 e_2 e_3 e_4}_{=\, B_1 \wedge B_2 \text{ (outer)}},\\
  A(B_1, B_2) &= \frac{1}{2}[B_1, B_2] = \underbrace{e_1 e_3}_{\in\, \Cl^2 \text{ (neither)}}.
\end{align*}
No binary inner/outer splitting of $B_1 B_2$ can capture this three-grade structure: the symmetric part $S(B_1, B_2)$ inextricably combines both the grade-$0$ inner contraction and the grade-$4$ outer product, while the antisymmetric part $A(B_1, B_2)$ lies purely in intermediate grade $2$.
\end{counterexample}

\begin{counterexample}[Failures of Metric Orthogonality --- Clause \textbf{(C4)}]\label{cex:orthog}
Clause \textbf{(C4)} also fails on two independent fronts:
\begin{enumerate}
  \item[(a)] \emph{Non-orthogonality of $S(X, Y)$ and $A(X, Y)$ even for positive-definite Euclidean $Q$:}
  In $\Cl_{2,0}(\mathbb{Q})$ ($Q(x_1, x_2) = x_1^2 + x_2^2$), the reversion metric $\langle Z, W \rangle_{\widetilde{Q}} = z_0 w_0 + z_1 w_1 + z_2 w_2 + z_{12} w_{12}$ is strictly positive definite on all of $\Cl_{2,0}(\mathbb{Q})$. Take the inhomogeneous elements
  \[
    X = 1 + e_1, \qquad Y = e_2 + e_1 e_2 = (1 + e_1)e_2.
  \]
  Since $(1 + e_1)^2 = 1 + 2 e_1 + e_1^2 = 2(1 + e_1)$ whereas $e_2(1 + e_1) = (1 - e_1)e_2$ and $(1 + e_1)(1 - e_1) = 1 - e_1^2 = 0$, we obtain
  \[
    X Y = 2(e_2 + e_1 e_2), \qquad Y X = (1 + e_1)(1 - e_1)e_2 = 0.
  \]
  Consequently,
  \[
    S(X, Y) = e_2 + e_1 e_2 = A(X, Y) \neq 0.
  \]
  Here $S(X, Y)$ and $A(X, Y)$ are \emph{nonzero and equal} (collinear rather than orthogonal), with
  \[
    \bigl\langle S(X, Y), A(X, Y) \bigr\rangle_{\widetilde{Q}} = 1^2 + 1^2 = 2 \neq 0, \qquad
    \langle X Y, X Y \rangle_{\widetilde{Q}} = 8 \neq 2 + 2 = \langle S, S \rangle_{\widetilde{Q}} + \langle A, A \rangle_{\widetilde{Q}}.
  \]
  \item[(b)] \emph{Isotropic and self-orthogonal blades for indefinite or degenerate quadratic forms (even on grade-$1$ vectors):}
  Definition~\ref{def:clifford} applies to \emph{arbitrary} quadratic forms $Q$.
  \begin{itemize}
    \item In split signature $\Cl_{1,1}(\mathbb{Q})$ with $Q(x_1, x_2) = x_1^2 - x_2^2$, let $u = e_1 + e_2$ ($Q(u) = 0$) and $v = e_1$. Then $u v = 1 - e_1 e_2 \neq 0$, with $S(u, v) = 1$ ($\langle S, S \rangle_{\widetilde{Q}} = 1$) and $A(u, v) = -e_1 e_2$ ($\langle A, A \rangle_{\widetilde{Q}} = -1$), so $\langle u v, u v \rangle_{\widetilde{Q}} = 1 + (-1) = 0$: a nonzero product has zero metric norm-square.
    \item In $\Cl_{1,1,1}(\mathbb{Q})$ with $Q(x_1, x_2, x_3) = x_1^2 - x_2^2 + x_3^2$, let $u' = e_1 + e_2$ and $v' = e_3$. Then $A' = A(u', v') = (e_1 + e_2)e_3 \neq 0$ satisfies $\langle A', A' \rangle_{\widetilde{Q}} = Q(u')Q(v') - B_Q(u', v')^2 = 0$ and $\langle A', A' \rangle_0 = 0$ (likewise $e_1 \wedge e_2 \neq 0$ in the exterior algebra $\Cl(\mathbb{Q}^2, 0)$). Thus $A' \in \mathrm{span}(A') \cap \mathrm{span}(A')^{\perp_{\widetilde{Q}}}$ is self-orthogonal, so the metric fails to induce a nondegenerate orthogonal direct sum.
  \end{itemize}
\end{enumerate}
\end{counterexample}

\begin{table}[htbp]
\centering
\footnotesize
\setlength{\tabcolsep}{3.8pt}
\renewcommand{\arraystretch}{1.22}
\begin{tabular}{@{}llllll@{}}
\toprule
\textbf{Algebra} & \textbf{Operands $(X, Y)$} & \textbf{Gr.} & $S(X, Y)$ & $A(X, Y)$ & \textbf{Status of (C2)--(C4)} \\
\midrule
$\Cl(V, Q)$ & $u, v \in V$ ($Q > 0$) & $(1, 1)$ & $B_Q(u, v)\cdot 1$ (inner) & $u \wedge v$ (outer) & \textbf{All True} (Thm.~\ref{thm:grade1}) \\
$\Cl_{2,0}(\mathbb{Q})$ & $(e_1,\, e_1 e_2)$ & $(1, 2)$ & $0$ & $e_2 = A \cdot B$ (inner) & \textbf{(C2), (C3) False} \\
$\Cl_{3,0}(\mathbb{Q})$ & $(e_1 e_2,\, e_3)$ & $(2, 1)$ & $e_{123} = X \wedge Y$ (outer) & $0$ & \textbf{(C2), (C3), (C4$'$) False} \\
$\Cl_{3,0}(\mathbb{Q})$ & $(e_1 e_2,\, e_2 e_3)$ & $(2, 2)$ & $0$ & $e_1 e_3 \in \Cl^2$ (Lie) & \textbf{(C3) False} ($X \wedge Y = 0$) \\
$\Cl_{4,0}(\mathbb{Q})$ & $(e_{12},\, e_{12}+e_{23}+e_{34})$ & $(2, 2)$ & $-1 + e_{1234}$ (both) & $e_{13} \in \Cl^2$ (neither) & \textbf{(C2), (C3) False} \\
$\Cl_{2,0}(\mathbb{Q})$ & $(1+e_1,\, e_2+e_{12})$ & mix & $e_2 + e_{12}$ & $e_2 + e_{12}$ ($= S$) & \textbf{(C4) False} ($\langle S, A \rangle_{\widetilde{Q}} = 2$) \\
$\Cl_{1,1,1}(\mathbb{Q})$ & $(e_1+e_2,\, e_3)$ & $(1, 1)$ & $0$ & $(e_1+e_2)e_3 \neq 0$ & $\langle A', A' \rangle_{\widetilde{Q}} = 0$ \\
\bottomrule
\end{tabular}
\caption{Summary of positive validity on grade-$1$ Euclidean vectors versus structural breakdowns across homogeneous multivectors, inhomogeneous elements, and indefinite metrics.}
\label{tab:summary}
\end{table}

\section{Formal Verification in Core Lean 4.19.0 and Exact Python}

All definitions, positive theorems, general parity laws, and counterexamples in this paper have been formalized from scratch in \textbf{Core Lean 4.19.0} (\texttt{lean/}, compiled with \texttt{-DwarningAsError=true}, zero external dependencies, zero \texttt{sorry}/\texttt{admit}/\texttt{native\_decide}/custom axioms) and cross-checked in exact rational Python (\texttt{verify.py}):
\begin{itemize}
  \item \texttt{CliffordDecomp/\allowbreak RatQ.lean}: Defines the typeclass \texttt{CommRing2 K} (commutative rings with invertible $2$) and constructs the exact field $\mathbb{Q}$ (\texttt{RatQ}) with full ring laws, field-inverse (\texttt{mul\_\allowbreak inv\_\allowbreak cancel}), positivity (\texttt{IsPos}), and reflective equality and inequality proofs.
  \item \texttt{CliffordDecomp/\allowbreak Algebra.lean}: Constructs the $2^n$-dimensional Clifford algebra \texttt{ClSpace K n} via the inductive Chevalley--Cayley--Dickson rule $(x_0 + e_0 x_1)(y_0 + e_0 y_1) = (x_0 y_0 + q_0 \widehat{x}_1 y_1) + e_0(x_1 y_0 + \widehat{x}_0 y_1)$, proving associativity (\texttt{mul\_\allowbreak assoc}), reversion anti-homomorphism (\texttt{rev\_\allowbreak mul}), the fundamental Clifford relation $v^2 = Q(v)\cdot 1$ (\texttt{vec\_\allowbreak sq\_\allowbreak eq\_\allowbreak quad}), and uniqueness of the universal Clifford lift (\texttt{lift\_\allowbreak unique}).
  \item \texttt{CliffordDecomp/\allowbreak GradingAndProducts.lean}: Formalizes grade projections (\texttt{gradeProj}), completeness $\sum_{k=0}^n \grade{X}{k} = X$ (\texttt{grade\_\allowbreak decomp\_\allowbreak complete}), Proposition~\ref{prop:c1} (\texttt{mul\_\allowbreak eq\_\allowbreak symPart\_\allowbreak add\_\allowbreak antisymPart}, \texttt{sym\_\allowbreak antisym\_\allowbreak decomp\_\allowbreak unique}), equivalence of \texttt{wedge} with top-grade projection (\texttt{homogeneous\_\allowbreak mul\_\allowbreak top\_\allowbreak grade}), contractions (\texttt{leftContract}, \texttt{rightContract}, \texttt{hestenesInner}), and positive definiteness of \texttt{metricForm} (\texttt{metricForm\_\allowbreak pos\_\allowbreak of\_\allowbreak ne\_\allowbreak zero}).
  \item \texttt{CliffordDecomp/\allowbreak PositiveTheorems.lean} and \texttt{Counterexamples.lean}: Establish the positive Grade-$1$ theorem (Theorem~\ref{thm:grade1}), the rational non-vacuity witness of Example~\ref{ex:nonvac}, the homogeneous grade-parity sign law (Theorem~\ref{thm:parity}), the universal vector--bivector parity reversal (Corollary~\ref{cor:vec-biv}), all rational refutations in Counterexamples~\ref{cex:cl2}--\ref{cex:orthog}, and the master disproof theorems (\texttt{clifford\_\allowbreak homogeneous\_\allowbreak decomp\_\allowbreak conjecture\_\allowbreak false} and \texttt{clifford\_\allowbreak product\_\allowbreak decomp\_\allowbreak conjecture\_\allowbreak false}).
\end{itemize}

\begin{thebibliography}{9}
\bibitem{Chevalley1954}
C.~Chevalley, \emph{The Algebraic Theory of Spinors}, Columbia University Press, New York, 1954.

\bibitem{HestenesSobczyk1984}
D.~Hestenes and G.~Sobczyk, \emph{Clifford Algebra to Geometric Calculus: A Unified Language for Mathematics and Physics}, D.~Reidel, Dordrecht, 1984.

\bibitem{Lounesto2001}
P.~Lounesto, \emph{Clifford Algebras and Spinors}, 2nd ed., London Mathematical Society Lecture Note Series, vol.~286, Cambridge University Press, Cambridge, 2001.

\bibitem{Dorst2007}
L.~Dorst, D.~Fontijne, and S.~Mann, \emph{Geometric Algebra for Computer Science}, Morgan Kaufmann, San Francisco, 2007.
\end{thebibliography}

\end{document}
